% !TEX root = ../../main.tex

\section{Towards Multiplicity One}\label{sec:multiplicity-one}

The sheet decomposition lemma (Lemma~\ref{lem:sheet-decomposition}) establishes integrality of the limit varifold $V$; we now refine it to a parity constraint on the density.

\begin{proposition}[Parity of the density function]\label{prop:density-parity}
Write $V = |\partial^*\Omega(x_0)| + V'$, where $V'$ is the non-negative varifold capturing the cancelling mass. Then at $\mathcal{H}^n$-a.e.\ point of $\spt\|V\|$:
\begin{enumerate}[label=\textup{(\roman*)}]
    \item at a pure cancellation point $p \in \spt\|V'\| \setminus \spt\|\partial^*\Omega(x_0)\|$, the density $\theta(p)$ is even;
    \item at a mixed regular point $p \in \reg(\partial^*\Omega(x_0)) \cap \spt\|V'\|$, the density $\theta(p)$ is odd.
\end{enumerate}
\end{proposition}

\begin{proof}
At a \emph{pure cancellation point} $p \in \spt\|V'\| \setminus \spt\|\partial^*\Omega(x_0)\|$, the set $\Omega(x_0)$ is locally constant near $p$. Say $B_r(p) \subset \Omega(x_0)$. If $x_i > x_0$, nestedness gives $B_r(p) \subset \Omega(x_i)$ and $\partial^*\Omega(x_i) \cap B_r(p) = \varnothing$, so only parameters $x_i < x_0$ contribute sheets. For such $x_i$, $\Omega(x_i) \subset \Omega(x_0)$, so $\Omega(x_i) \cap B_r(p) = B_r(p)$ minus thin slabs, each bounded by two sheets. Therefore $k$ is even. (The case $B_r(p) \subset M \setminus \overline{\Omega(x_0)}$ is analogous.) The geometric mechanism is \emph{folding}: each pair of sheets bounds a slab that collapses as $x_i \to x_0$, contributing multiplicity~$2$.

At a \emph{mixed regular point} $p \in \reg(\partial^*\Omega(x_0)) \cap \spt\|V'\|$, the boundary $\partial^*\Omega(x_0)$ is smooth near $p$\xinze{Was: cited removed Federer Proposition; smooth by def of $\reg$}, and the $k$ sheets from Lemma~\ref{lem:sheet-decomposition} are ordered by height: $u_i^1 < \cdots < u_i^k$. Since $\Omega(x_i)$ occupies alternating slabs and $\Omega(x_i) \to \Omega(x_0)$ in $L^1$, exactly one sheet (the \emph{tracking sheet}) satisfies $\|u_i^{j_0} - u_0\|_{L^\infty} \to 0$, where $u_0$ parametrizes $\partial^*\Omega(x_0)$. The remaining $k - 1$ cancelling sheets come in pairs by the same folding argument. Therefore $k - 1$ is even, i.e., $k$ is odd.

For $n \leq 6$, the Federer dimension bound~\cite[Section~37]{Sim84}\xinze{Was: Federer regularity Proposition} gives $\partial\Omega(x_0) = \partial^*\Omega(x_0)$, so there are no singular points to consider.
\end{proof}

Thus multiplicity one ($\theta \equiv 1$) is equivalent to $V' = 0$ everywhere, which is precisely the no-mass-cancellation condition (Definition~\ref{def:mass-cancellation}).

A natural strategy for proving multiplicity one is the following. Suppose $V = k|\Sigma|$ with $k \geq 2$ and $\Sigma$ a smooth closed embedded minimal hypersurface (as guaranteed by Theorem~\ref{thm:non-excessive-main}). One attempts to construct an $I$-replacement family (Definition~\ref{def:excessive-interval}) on an interval $I \ni x_0$ by performing surgery on near-maximal slices: for $x$ with $\mathrm{Per}(\Omega(x))$ close to $W = k \cdot \mathcal{H}^n(\Sigma)$, replace $\Omega(x)$ inside a tubular neighborhood $B_\varepsilon(\Sigma)$ by $\Omega(x_0)$, removing the $k - 1$ cancelling sheet pairs. Since $\mathrm{Per}(\Omega(x_0)) \leq \mathcal{H}^n(\Sigma) = W/k$, and the cutting cost on $\partial B_\varepsilon(\Sigma)$ vanishes as $x \to x_0$ by $L^1$ convergence, the modified perimeter satisfies $\mathrm{Per}(\tilde{\Omega}(x)) \approx W/k \ll W$. This would make $x_0$ excessive, contradicting non-excessiveness.

The obstruction is that $L^1$ convergence of Caccioppoli sets does not control the spatial distribution of perimeter. For the specific min-max sequence $x_i \to x_0$, varifold convergence $|\partial^*\Omega(x_i)| \to k|\Sigma|$ guarantees that the perimeter of $\Omega(x_i)$ concentrates in $B_\varepsilon(\Sigma)$, so the surgery reduces perimeter as claimed. But for a general parameter $x$ near $x_0$, the boundary $\partial^*\Omega(x)$ need not be near $\Sigma$: a slice with $\mathrm{Per}(\Omega(x)) \approx (k-1)\mathcal{H}^n(\Sigma)$ and $\mathrm{Per}(\Omega(x), B_\varepsilon(\Sigma)) \approx 0$ would see its perimeter \emph{increase} to $\approx (k-1)\mathcal{H}^n(\Sigma) + \mathcal{H}^n(\Sigma) = W$ after the replacement, violating the mass bound required of an $I$-replacement family. The $I$-replacement definition demands $\limsup$ mass $< W$ for \emph{all} $x \in I$, and one cannot restrict the surgery to the subsequence where it is effective.

This gap reflects a general difficulty: the non-excessive framework controls the limit varifold through one-sided homotopic minimizing properties of critical slices, but does not directly constrain the geometry of nearby non-critical slices. Zhou's proof of multiplicity one~\cite{Zho20} overcomes analogous obstructions using the substantially deeper machinery of prescribed mean curvature (PMC) min-max theory, which provides finer control on the interaction between the sweepout parameter and the geometry of individual slices.

\begin{question}\label{q:mult-one}
Can multiplicity one be established within the CLS non-excessive framework, using nestedness and the sheet decomposition (Lemma~\ref{lem:sheet-decomposition}) to rule out mass cancellation?
\end{question}
